\section{도구 선택 성능}\label{sec:exp:performance}

본 절은 중심 주장 1 — 회전 기반 스티어링이 멀티-툴 도구 선택을 개선한다 — 을 지지하는 증거를 쌓는다. 동일한 프롬프트, 스코러, greedy 디코딩 아래에서 여섯 방법을 비교한다: \texttt{no\_steer}, 우리 정지 K 구현 \texttt{ocq\_bias\_a0.3}, 외부 canonical SEKA, canonical AdaSEKA, 그리고 본 연구의 Q-커버리지(\texttt{ocq\_qbias\_b-0.03})와 레이어 적응형 K+Q(\texttt{ocq\_ladapt\_k0.05\_q-0.03}).

\subsection{MetaTool Subtask4}\label{sec:exp:metatool}

MetaTool Subtask4(N=497)는 각 질의에 정확히 두 도구의 정답이 지정되는, 순차 coverage를 가장 깨끗하게 측정하는 벤치마크다. 표~\ref{tab:e1-subtask4}의 잠금된 수치는 이미 핵심 이야기를 드러낸다. 같은 카탈로그 온톨로지 기저 위에서 정지 key-side 연산자는 $-4.57$pp의 심각한 음의 부호를 보이는 반면, 부호만 뒤집어 query에서 빼는 본 연구의 두 연산자는 각각 $+2.28$pp(Q-커버리지)와 $+2.08$pp(레이어 적응형)의 양의 이득을 낸다. 외부 canonical SEKA/AdaSEKA 수치는 부록~\ref{app:seka-repro}의 재현 gate를 통과한 뒤 채워진다.

\begin{table}[t]
\centering
\caption{MetaTool Subtask4 (N=497), Qwen2.5-7B-Instruct. 굵은 값 locked, 이탤릭은 예상 방향.}
\label{tab:e1-subtask4}
\small
\begin{tabularx}{\linewidth}{@{}Xrrr@{}}
\toprule
방법 & F1 & Exact & $\Delta$F1 \\
\midrule
\texttt{no\_steer} & \textbf{0.7307} & \textbf{0.5252} & --- \\
canonical SEKA amp=1.0 & \emph{placeholder} & \emph{placeholder} & $\ll 0$ \\
canonical AdaSEKA amp=1.0 & \emph{placeholder} & \emph{placeholder} & $\ll 0$ \\
정지 K \texttt{a0.3} & \textbf{0.6850} & \textbf{0.4728} & $\mathbf{-4.57}$ \\
Q-커버리지 (ours) & \textbf{0.7535} & \textbf{0.5332} & $\mathbf{+2.28}$ \\
레이어 적응형 K+Q (ours) & \textbf{0.7514} & \textbf{0.5473} & $\mathbf{+2.08}$ \\
\bottomrule
\end{tabularx}
\end{table}

\subsection{$\tau^2$-bench retail}\label{sec:exp:tau2}

독립 도메인 검증을 위해 $\tau^2$-bench retail의 첫 턴 action-set selection을 평가한다. 정답 도구 수가 4에서 13까지 분포하는 이 체제는 MetaTool의 $k=2$와 대비되며, 이론의 regime-split 예측을 시험할 수 있는 환경을 제공한다. 표~\ref{tab:e2-tau2}의 N=30 smoke는 이미 예상 패턴을 보여준다. Q-커버리지는 Qwen F1을 $+3.74$pp 끌어올리는 반면 레이어 적응형은 $-3.73$pp 떨어진다. 이 비대칭은 우연이 아니라 따름정리~\ref{cor:regime}의 정량 예측 — K-imprinting의 한계 이득이 긴 시퀀스에서 감쇠 — 과 정합적이다. 같은 연산자가 $k=2$에서는 최고이고 $k\ge 4$에서는 열세라는 패턴은 우리 이론이 이미 예고한 것이다. N=114 full이 통계적 확정을 제공할 것이다.

\begin{table}[t]
\centering
\caption{$\tau^2$-bench retail. Qwen2.5-7B. N=30 smoke locked.}
\label{tab:e2-tau2}
\small
\begin{tabularx}{\linewidth}{@{}lXrrr@{}}
\toprule
N & 방법 & F1 & Recall & $\Delta$F1 \\
\midrule
\multirow{5}{*}{30 (smoke)}
 & \texttt{no\_steer} & \textbf{0.5735} & \textbf{0.5340} & --- \\
 & Q-커버리지 (ours) & \textbf{0.6109} & \textbf{0.5717} & $\mathbf{+3.74}$ \\
 & 레이어 적응형 K+Q (ours) & \textbf{0.5362} & \textbf{0.4831} & $\mathbf{-3.73}$ \\
 & canonical SEKA amp=1.0 & \emph{placeholder} & \emph{placeholder} & $\le 0$ \\
 & canonical AdaSEKA amp=1.0 & \emph{placeholder} & \emph{placeholder} & $\le 0$ \\
\midrule
\multirow{3}{*}{114 (full)}
 & \texttt{no\_steer} & \emph{placeholder} & \emph{placeholder} & --- \\
 & Q-커버리지 (ours) & \emph{placeholder} & \emph{placeholder} & $> 0$ \\
 & 레이어 적응형 K+Q (ours) & \emph{placeholder} & \emph{placeholder} & $\approx 0$ 또는 $<0$ \\
\bottomrule
\end{tabularx}
\end{table}

\subsection{승리의 원인: 포지션별 분해}\label{sec:exp:mechanism}

집계 F1이 예측 방향과 일치한다는 것은 필요조건일 뿐이다. \emph{왜} 이득이 생기는지에 대한 진짜 증거는 분해에서 나온다. 정리~\ref{thm:no-memory}의 관측 가능한 귀결은 ``정지 key-side가 첫 도구를 두 번째 단계에서 다시 고르게 유도한다''이므로, 데이터의 특정 위치 — 두 번째 포지션 — 에 특정한 흔적이 있어야 한다. \texttt{repeated\_first\_tool\_rate}이 SEKA에서 \emph{올라가고}, 우리 연산자에서 \emph{내려가야} 하며, 이와 동시에 \texttt{second\_distinct\_hit\_rate}이 우리 연산자에서 \emph{올라가야} 한다. 이 세 관측이 동시에 성립할 때만 이득이 ``반복을 지운 결과''라는 해석이 지지된다.

그림~\ref{fig:stepwise}와 표~\ref{tab:e3-mechanism}이 그 패턴을 보여준다. 네 방법에 대해 네 포지션 지표를 나란히 놓으면 정지 K 계열 모두에서 \texttt{repeated\_first\_tool\_rate}이 \texttt{no\_steer}보다 뚜렷이 높아지고 \texttt{second\_distinct}가 가장 많이 떨어진다. 우리 연산자는 반복을 약화시키면서 두 번째 포지션의 정확도를 끌어올린다. 이 네 수치의 \emph{동시 방향성}이 본 논문의 핵심 증거다. F1 하나였다면 ``같은 기저를 쓰는 두 연산자가 우연히 다르게 나왔다''는 해석이 가능했겠지만, 네 수치가 같은 이야기를 가리킬 때는 그런 우연 해석이 어렵다.

\begin{figure}[t]
\centering
\includegraphics[width=0.88\linewidth]{fig3_stepwise.pdf}
\caption{포지션별 hit rate. 정지 K-계열의 \texttt{repeated\_first} 상승과 우리 연산자의 감소가 동시에 관측되는 것이 따름정리~\ref{cor:seka-repeat}의 직접 확증이다.}\label{fig:stepwise}
\end{figure}

\begin{table}[t]
\centering
\caption{MetaTool Subtask4 N=497의 포지션별 hit rate. 화살표는 정리가 예측하는 방향.}
\label{tab:e3-mechanism}
\small
\begin{tabularx}{\linewidth}{@{}Xcccc@{}}
\toprule
방법 & first\_hit & second\_hit & second\_distinct & repeated\_first $\downarrow$ \\
\midrule
\texttt{no\_steer} & \emph{baseline} & \emph{baseline} & \emph{baseline} & \emph{baseline} \\
canonical SEKA amp=1.0 & $\approx$ & $\downarrow\downarrow$ & $\downarrow\downarrow$ & $\boldsymbol{\uparrow\uparrow}$ \\
정지 K \texttt{a0.3} & $\approx$ & $\downarrow\downarrow$ & $\downarrow\downarrow$ & $\boldsymbol{\uparrow\uparrow}$ \\
Q-커버리지 (ours) & $\approx$ & $\uparrow$ & $\uparrow$ & $\downarrow$ \\
레이어 적응형 K+Q (ours) & $\boldsymbol{\uparrow}$ & $\boldsymbol{\uparrow\uparrow}$ & $\boldsymbol{\uparrow\uparrow}$ & $\downarrow$ \\
\bottomrule
\end{tabularx}
\end{table}

\subsection{ablation: 대안적 해석을 하나씩 반박}\label{sec:exp:ablations}

중심 주장을 위협할 수 있는 세 가지 대안 해석이 있다. 첫째, 이득이 기저 품질 덕분이라는 해석. 둘째, 효과가 단지 ``약하게 쓰니까 안전했다''는 해석. 셋째, 연산자가 일종의 ``확신 증가제''로 작동해 많이 쓸수록 좋다는 해석. 본 절은 이 셋을 각각 반박하는 ablation을 제시하고, 마지막으로 정리~\ref{thm:bound}의 진단 bound가 실제로 정보적임을 경험적으로 확인한다.

\paragraph{기저 특이성.}
표~\ref{tab:e4-basis}는 ``기저 덕분'' 해석을 분리한다. 정지 K는 feature-shuffled와 random 기저에서 완전히 0으로 붕괴한다. 같은 연산자가 카탈로그 기저에서만 $-4.57$pp의 ``온전한 음의 효과''를 유지한다는 사실은 기저 방향이 정지 K의 행동을 전적으로 결정한다는 뜻이다. Q-커버리지는 기저가 바뀌어도 F1을 0.7068 이상 유지하며 상대적으로 덜 민감하다. 이 비대칭은 두 연산자가 부분공간 기하에 다른 의존도를 가진다는 의미다. PCA-of-K 기저의 결과가 마지막 단계를 결정한다. PCA가 real과 구별되지 않으면 ``카탈로그 온톨로지가 필수''라는 강한 주장이 약화되고, PCA가 중간값 또는 neutral이면 ``저랭크 구조가 필요하되 카탈로그가 최적''이라는 약한 주장으로 좁혀진다. 어느 경우든 random과 shuffled의 붕괴는 변하지 않으므로 ``아무 방향이든 된다''는 대안은 이미 배제된다.

\begin{table}[t]
\centering
\caption{Qwen2.5-7B, Subtask4 N=497: 기저 $\times$ 연산자. 굵은 값 locked.}
\label{tab:e4-basis}
\small
\begin{tabularx}{0.9\linewidth}{@{}Xcc@{}}
\toprule
기저 & Q-커버리지 F1 & 정지 K F1 \\
\midrule
Real $B_{\mathrm{ont}}$ (ours) & \textbf{0.7471} & \textbf{0.6850} \\
Feature-shuffled & \textbf{0.7254} & \textbf{0.0000} \\
Random orthonormal & \textbf{0.7068} & \textbf{0.0000} \\
PCA-of-K (calibration) & \emph{placeholder} & \emph{placeholder} \\
\bottomrule
\end{tabularx}
\end{table}

\paragraph{``약해서 안전했다''의 반박.}
자연스러운 반론은 real 기저가 null control보다 교란을 덜 주기 때문에 안전한 것 아니냐는 것이다. 표~\ref{tab:e-perturb-mag}가 이 설명을 직접 반박한다. real 온톨로지 기저는 실제 attention perturbation magnitude(Qwen L=13, $\alpha=0.3$, N=30)에서 random과 feature-shuffled 기저보다 \emph{더 큰} 교란을 만든다 — 평균 $621.30$ vs random $399.56$ vs feature-shuffled $291.98$. 그럼에도 K-side 비교에서 유일하게 안정적인 방향은 real 기저뿐이다. 곧 real 기저가 ``이기는 이유는 약해서가 아니라 특정한 안정 부분공간에 정렬되어 있기 때문''이다.

\begin{table}[t]
\centering
\caption{Qwen L=13, $\alpha=0.3$, N=30에서의 effective attention perturbation magnitude. Real 기저가 가장 크게 교란하는데도 안정성은 real에서만 유지된다.}
\label{tab:e-perturb-mag}
\small
\begin{tabularx}{0.9\linewidth}{@{}Xccc@{}}
\toprule
기저 & 평균 & 중앙값 & 표준편차 \\
\midrule
Real $B_{\mathrm{ont}}$ (ours) & \textbf{621.30} & \textbf{593.71} & 317.59 \\
Random orthonormal & 399.56 & 360.01 & 235.86 \\
Feature-shuffled & 291.98 & 241.87 & 169.97 \\
\bottomrule
\end{tabularx}
\end{table}

\paragraph{``많이 쓸수록 좋다''의 반박.}
세 번째 반론은 Q-커버리지가 단지 일종의 confidence boost로 작동해서 $|\beta|$가 클수록 좋아져야 한다는 것이다. 이 해석도 데이터가 배제한다. 그림~\ref{fig:beta-sweep}의 Q-only $\beta$ sweep은 좁은 작동 구간을 드러낸다. 평가한 값들 중 $\beta = -0.1$ 근처가 최고이며, $\beta = -0.7$에 이르면 F1이 0으로 완전 붕괴한다. 패턴은 ``bias가 클수록 좋다''가 아니라 ``작은 음의 Q-side projection은 도움이 되지만, 그 밖의 대부분 설정은 오히려 성능을 깎는다''에 가깝다. 이 좁은 구간은 조잡한 confidence 증가보다 coverage effect 해석과 잘 맞는다.

\begin{figure}[t]
\centering
\includegraphics[width=0.75\linewidth]{fig_qbias_ablations.pdf}
\caption{Qwen Subtask4 Q-커버리지 $\beta$ sweep. 좁은 유효 구간은 confidence-boost 해석과 어긋나고 coverage 해석과 부합한다.}\label{fig:beta-sweep}
\end{figure}

\paragraph{진단 bound의 경험적 유효성.}
마지막으로 정리~\ref{thm:bound}의 부등식이 공허하지 않은지 확인한다. Qwen2.5-7B(L=13)에서 $\alpha=0.3$의 정지 K 개입에 대해 N=100, 전 head(2800개 head-query 조합)에서 측정한 LHS/RHS 비율은 중앙값 $2.36\times 10^{-8}$, pass rate 100\%로 bound는 매우 보수적이지만 안정적으로 성립한다. 이것은 Q-side와 K-side 개입을 공통 척도로 비교할 때 정리가 실제로 정보적 진단이 될 수 있음을 뒷받침한다.

\paragraph{단일 도구 체제에서의 거동.}
마지막 보조 증거는 반대 방향에서 주장을 강화한다. MetaTool Subtask1의 단일 도구 체제에서 Qwen에 정지 K-bias($\alpha=0.3$)를 가하면 $+11.16$pp의 \emph{큰 양의 이득}을 준다. 언뜻 본 논문의 주장과 어긋나 보이지만 오히려 확증한다. 도구가 하나뿐인 체제에서는 ``첫 선택을 반복하라''는 편향이 중립 또는 도움이 되며, 이것이 SEKA 계열이 단일 개념 편집에서 강한 이유와 일치한다. 본 연구의 두 연산자는 SEKA의 범용 대체가 아니라 멀티-툴 과제에 특화된 설계다.

\section{효율성}\label{sec:exp:efficiency}

중심 주장 2 — 회전 기반 스티어링이 대안 대비 효율적이다 — 는 실험이 아니라 네 축의 직접 비교로 지지된다. 표~\ref{tab:e-efficiency}가 본 방법을 tool-aware 파인튜닝, 프롬프트 스캐폴딩, KV-cache 최적화, 정지 활성 스티어링과 비교한다.

\begin{table}[t]
\centering
\caption{효율성 비교. 본 방법은 추론 시 hook으로만 작동하여 학습과 토큰 오버헤드가 0이며, 저장은 per-head 기저만 추가로 든다(Qwen 7B 기준 $\approx 2.6$ MB).}
\label{tab:e-efficiency}
\small
\begin{tabularx}{\linewidth}{@{}Xrrrr@{}}
\toprule
접근 & 학습 비용 & 토큰 오버헤드 & 추론 지연 & 저장 \\
\midrule
Tool-aware 파인튜닝 (Gorilla, ToolLlama) & $\mathcal{O}(\text{일})$ & 0 & 0 & 전체 모델 \\
프롬프트 스캐폴딩 (ReAct) & 0 & $+200\sim 500$ 토큰 & 프롬프트 길이 비례 & 0 \\
KV-cache 최적화 (KIVI) & 0 & 0 & 근소 감소 & 캐시 $-30\%$ \\
정지 활성 스티어링 (SEKA) & 사영 학습($\mathcal{O}$(분)) & 0 & 근소 증가 & 사영 행렬 \\
\textbf{본 논문 (Q-커버리지, K+Q)} & \textbf{0} & \textbf{0} & \emph{측정 예정} & \textbf{$\approx 2.6$ MB} \\
\bottomrule
\end{tabularx}
\end{table}

``측정 예정'' 칸은 두 보완 측정으로 채워진다. 추론 지연 측정은 네 조건(\texttt{no\_steer}, Q-커버리지, 레이어 적응형 K+Q, canonical SEKA) 아래에서 per-query wall-clock을 기록하고, 메모리 프로파일은 hook 설치 전후의 peak GPU memory로 실제 추가 footprint를 확인한다. hook이 attention에 $O(Tdr)$ 행렬 곱만 추가하므로 이론적 오버헤드는 전체 추론 시간의 1\% 미만일 것이 예상되며, 측정이 이 예측을 수치로 확정한다. 특기할 점은 KV-cache 최적화 계열과 본 방법이 서로 다른 목표(메모리 vs 과제 성능)를 겨냥하며 \emph{결합 가능}하다는 것이다. KIVI 등이 캐시를 압축하는 동안 본 방법은 같은 캐시 위에서 과제 수준의 선택을 개선한다.
